\paragraph{RQ1: Construction and capacity.}
How can JHQ's hierarchy be trained, encoded and arranged as a searchable
GPU-resident index, and what construction cost and tested capacity result?
The evidence is the physical layout, train/add records and resident-memory
account (\cref{sec:7-7}). An unsuccessful build or an allocation beyond the
device budget would limit the capacity claim; a resident sample alone does
not establish construction peak.

\paragraph{RQ2: Correct and efficient execution.}
Which representation-aware scoring, retention and scheduling choices
improve GPU search while respecting the specified score and tie semantics?
Paired LUT, packed-load, refinement and scan experiments connect to reference
checks and batch sensitivity (\cref{sec:7-3,sec:7-4,sec:7-6}). A non-tie
selection error would challenge the correctness contract; consistent
slowdowns would challenge the proposed execution choice.

\paragraph{RQ3: Budget quality and total cost.}
Can unlabelled workload agreement choose a useful survivor budget, and under
what sampling and reuse conditions is calibration worthwhile? Disjoint query
evaluation compares fixed budgets, sampled/FULL rules, enumeration and a
GT-labelled oracle (\cref{sec:7-5}). Loss crossings challenge reliability;
no positive saving after calibration challenges the net-cost benefit.
The oracle is an evaluation control, not deployable label-free logic.

\paragraph{RQ4: Competitive operating regions.}
At which measured recall targets does the complete system compete with
modern GPU ANN alternatives under a common host-to-host boundary?
The answer uses independently tuned E1 frontiers and explicit memory
definitions (\cref{sec:7-2}). Int8 graph competitors, high-recall reversals
and absent coverage are part of the answer. The question does not seek a
universal ranking across hardware, memory budgets or request sizes.

Chapter 3 turns these questions into controlled comparisons. Together they
separate what was built, whether its execution is correct, what effort a
workload requires, and whether the resulting system has useful operating
regions.
